\documentclass[11pt]{article}
\usepackage[a4paper,margin=26mm]{geometry}
\usepackage{amsmath,amssymb,amsthm,mathtools,mathrsfs,hyperref}
\hypersetup{hypertexnames=false}
\setlength{\emergencystretch}{2em}
\newtheorem{theorem}{Theorem}
\newtheorem{lemma}[theorem]{Lemma}
\newcommand{\C}{\mathbb C}
\newcommand{\R}{\mathbb R}
\newcommand{\HH}{\mathcal H}
\newcommand{\PP}{\mathcal P}
\newcommand{\V}{\mathcal V}
\newcommand{\diag}{\operatorname{diag}}
\title{The complete complex amplitude in gamma sum descent}
\author{}
\date{13 September 2026}
\begin{document}
\maketitle

\section{Original objects and the question answered}
The delivered Gamma Convolution Descent note constructs the conditional
expectation of the squared arithmetic amplitude. We construct the exact
comparison for its full complex amplitude. The two resulting source
Grams add to the original arithmetic Gram. We prove a strict positive
contribution of the retained relative component for every full quartet,
and give its exact map into the original minimum quotient metric.
All measures below retain their masses. The variable is the original
sum $u=t_1+\cdots+t_k$, and the arithmetic variable remains $S=k/2+iu$.

Fix $\lambda>0$ and an integer $k\geq1$. Retain
\begin{equation}
 r_\lambda(t)=\frac{|\Gamma(\lambda+it/2)|^2}{2\pi},\quad
 c_\lambda=2^{1-2\lambda}\Gamma(2\lambda),\quad
 r_{\lambda,k}=r_\lambda^{*k}=
 \frac{c_\lambda^k}{c_{k\lambda}}r_{k\lambda}.
 \tag{GP.1}
\end{equation}
Thus $\int r_{\lambda,k}=c_\lambda^k$. The beta-integral proof of
this identity, with its original Jacobian $t=2v$, is retained in
the cumulative reader, TG.7--TG.8, and the delivered note, (5)--(9).
In particular
\begin{equation}
 \int e^{\theta u}r_{\lambda,k}(u)\,du
       =c_\lambda^k(\cos\theta)^{-2k\lambda},
       \qquad |\theta|<\pi/2.
 \tag{GP.2}
\end{equation}
Let $h$ be the original finite polynomial packet of complete orders of
actual zeros of $g=2\xi$, and include $h=1$ when considering the analytic
seed. Define the entire quotient $v_h=g/h$ and keep
\begin{equation}
 A(t)=\frac{v_h(1/2+it)}{\Gamma(\lambda+it/2)},\quad
 B(t)=|A(t)|^2,\quad w(t)=B(t)r_\lambda(t),\quad
 \mu_h=\int w(t)\,dt.
 \tag{GP.3}
\end{equation}
Use one of the references proved in the delivered note for which
$B\leq C_h<\infty$. In particular a full quartet of degree $d\geq4$
uses its original $\lambda=1/4$. The seed uses $\lambda=13/4$ and
$C_1=1024/\sqrt\pi$. The statements involving decay explicitly use
the quartet reference. The function $A$ is nonzero and analytic on
the real line: $v_h$ is a nonzero entire function and reciprocal gamma
is entire and nonzero on this line. Its real zeros form a discrete
set. No value at a removed packet zero is assigned using an uncancelled
fraction.

Put $d\nu(\mathbf t)=\prod_i r_\lambda(t_i)\,d\mathbf t$ and
$d\nu_\Sigma(u)=r_{\lambda,k}(u)du$. In
$\HH=L^2(\nu)$ and $\HH_\Sigma=L^2(\nu_\Sigma)$ write
\begin{equation}
 \begin{gathered}
 Uf(\mathbf t)=f(\textstyle\sum_i t_i),\quad C=U^*,\quad P=UC,
 \quad a(\mathbf t)=\prod_i A(t_i),\\
 b_k=C(|a|^2),\quad a_k=C(a),\quad d_k=b_k-|a_k|^2.
 \end{gathered}
 \tag{GP.4}
\end{equation}
For $k=1$, $C=U=I$, $a_1=A$, $b_1=B$, $d_1=0$. For $k\geq2$
the conditional integral $CF$ is precisely
\begin{equation}
 (CF)(u)=\frac{1}{r_{\lambda,k}(u)}
 \int_{\R^{k-1}}F(t_1,\ldots,t_{k-1},u-\textstyle\sum_{i<k}t_i)
 \prod_{i<k}r_\lambda(t_i)\,r_\lambda(u-\textstyle\sum_{i<k}t_i)
 \,d^{k-1}t.
 \tag{GP.5}
\end{equation}
These are probability measures on individual fibres, derived by dividing
by their stated fibre mass; neither original measure is replaced by a
probability measure. Fubini proves $CU=I$ and the isometry of $U$.
Jensen's inequality in (GP.5) proves $d_k\geq0$ and
$0\leq b_k\leq C_h^k$. For $k\geq2$, $b_k(u)>0$ at every real $u$:
the discrete zeros of each factor exclude only countably many
hyperplanes in this fibre, which have zero Lebesgue measure.

\section{Unitary transport of the actual conditional projection}
\begin{theorem}[Complete amplitude comparison]
Use the actual arithmetic spaces
\[
 \HH^h=L^2(\R^k,|a|^2d\nu),\qquad
 \HH^h_\Sigma=L^2(\R,b_kd\nu_\Sigma).
\]
The maps
\begin{equation}
 T:\HH^h\longrightarrow\HH,\quad F\longmapsto aF,
 \qquad T_\Sigma:\HH^h_\Sigma\longrightarrow\HH_\Sigma,
       \quad f\longmapsto\sqrt{b_k}f
 \tag{GP.6}
\end{equation}
are surjective isometries. The positive square root in $T_\Sigma$
is a specified comparison map; the original amplitude in $T$ remains
complex. The actual arithmetic sum inclusion $U_hf=f(\sum t_i)$
and its adjoint $C_h^{\rm sum}$ satisfy
\begin{equation}
 C_h^{\rm sum}F=\frac{C(|a|^2F)}{b_k},\qquad
 J=TU_hT_\Sigma^{-1}=M_aUM_{b_k^{-1/2}},\quad
 J^*=M_{b_k^{-1/2}}CM_{\bar a},\quad J^*J=I.
 \tag{GP.7}
\end{equation}
All quotients here are defined almost everywhere. In particular their
operator extensions are bounded even when the individual inverse
multiplier is unbounded. The transported arithmetic projection is
\begin{equation}
 Q=TP_hT^{-1}=JJ^*,\qquad
 QF=a\,U\left(\frac{C(\bar a F)}{b_k}\right),
 \quad P_h=U_hC_h^{\rm sum}.
 \tag{GP.8}
\end{equation}
Thus $P$ and $Q$ are related by the displayed exact transport from
the original arithmetic spaces.
\end{theorem}
\begin{proof}
The squared norms in (GP.6) agree by direct substitution. The zero set
of $a$ is a union of null hyperplanes; $b_k$ is positive almost
everywhere, including $k=1$. Division off those null sets constructs
the two inverses with their exact squared norms. Applying (GP.5) to
the arithmetic fibre density gives $C_h^{\rm sum}$. Pairing
$M_aUM_{b_k^{-1/2}}f$ with $F$ proves the stated adjoint. The identity
$C(|a|^2Uf)=b_kf$ then gives $J^*J=I$. Conjugating $P_h$ gives
(GP.8). These norm identities first hold for bounded functions
supported on finite-measure sets on which the denominators are
bounded below; truncation and the isometries extend them to the
indicated complete Hilbert spaces. This also proves that no unbounded
multiplier composition is being asserted on a larger domain.
\end{proof}

\begin{theorem}[Exact fibre angle and retained defect]
Define $\gamma_k=a_k/\sqrt{b_k}$ almost everywhere. Then
\begin{equation}
 U^*J=M_{\gamma_k},\quad
 J=UM_{\gamma_k}+N,\quad U^*N=0,\quad
 N^*N=M_{1-|\gamma_k|^2}=M_{d_k/b_k}.
 \tag{GP.9}
\end{equation}
For every real $u$ at which $b_k>0$, the two fibre projections onto
$1$ and $a/\sqrt{b_k}$ obey
\begin{equation}
 \operatorname{Tr}_{u}(P_u-Q_u)^2=2(1-|\gamma_k(u)|^2),\qquad
 \|[P_u,Q_u]\|_{\rm HS,u}^2
       =2|\gamma_k(u)|^2(1-|\gamma_k(u)|^2).
 \tag{GP.10}
\end{equation}
These are finite-rank traces in the specified fibre Hilbert space,
not a claim that the global decomposable operators are trace class.
The complete unscaled variance also has the formula
\begin{equation}
 d_k(u)=\frac12\iint|a(\mathbf t)-a(\mathbf s)|^2
                       \,d\nu_u(\mathbf t)d\nu_u(\mathbf s),
 \tag{GP.11}
\end{equation}
where both $\nu_u$ are the exact conditional measures in (GP.5).
\end{theorem}
\begin{proof}
Applying $C$ to $Jf$ gives $\gamma_k f$. Orthogonal decomposition
with $P$ gives (GP.9), and taking squared norms uses $J^*J=I$.
On a fibre let $e_0=1$ and $v=a/\sqrt{b_k}$, so both have norm
one and $\langle e_0,v\rangle=\gamma_k$, using the inner product
antilinear in its first argument. If $|\gamma_k|<1$, put
$e_1=(v-\gamma_ke_0)/\sqrt{1-|\gamma_k|^2}$. In the displayed
orthonormal comparison frame, without altering either original
vector, the two projections are
\[
 P_u=\begin{pmatrix}1&0\\0&0\end{pmatrix},\qquad
 Q_u=\begin{pmatrix}
 |\gamma_k|^2&\gamma_k\sqrt{1-|\gamma_k|^2}\\
 \overline{\gamma_k}\sqrt{1-|\gamma_k|^2}&1-|\gamma_k|^2
 \end{pmatrix}.
\]
Matrix multiplication proves (GP.10); both vanish when
$|\gamma_k|=1$. Expanding the square in (GP.11), and using the
unit total mass of each conditional measure, gives
$C(|a|^2)-|Ca|^2$. No phase is discarded in this calculation.
\end{proof}

\section{Strictness for an actual quartet, with its exact constant}
For a full quartet retain $\lambda=1/4$, degree $d\geq4$ and
$T_h=\max(1,2\max_\rho|\rho-1/2|)$. The delivered proof gives
\begin{equation}
 |A(t)|^2\leq \frac{25\,2^{2d}}{\sqrt\pi}|t|^{6-2d}
                  \quad (|t|\geq T_h).
 \tag{GP.12}
\end{equation}
Here is its direct calculation. In the unchanged amplitude,
\[
 A(t)=-\pi^{-1/4}(t^2+1/4)e^{-it\log\pi/2}
                    \zeta(1/2+it)/h(1/2+it).
\]
The continuation formula
$\zeta(s)=s/(s-1)-s\int_1^\infty\{x\}x^{-s-1}dx$
gives $|\zeta(1/2+it)|\leq1+2\sqrt{t^2+1/4}
\leq2(1+|t|)$. For $|t|\geq T_h$, every retained root factor has
modulus at least $|t|/2$, so $|h|\geq(|t|/2)^d$.
Also $(t^2+1/4)^2\leq25t^4/16$ and
$4(1+|t|)^2\leq16t^2$ for $|t|\geq1$. Multiplication proves
(GP.12). Consequently $A(t)\to0$ at both real ends.

\begin{theorem}[Strict relative contribution on every sum fibre]
For every such actual quartet, every $k\geq2$, and every real $u$,
\begin{equation}
 d_k(u)>0,\qquad |\gamma_k(u)|<1.
 \tag{GP.13}
\end{equation}
Thus the relative component is nonzero at every fibre. It is given
exactly by (GP.9) and (GP.11); strictness does not assign a uniform
positive lower bound to it.
\end{theorem}
\begin{proof}
Fix $u$. If $d_k(u)=0$, (GP.11), or equality in the variance
identity, says that the continuous function $a$ is constant almost
everywhere on the fibre. Its conditional density is strictly positive
everywhere, so continuity makes it constant everywhere. For $k=2$,
this would make $A(t)A(u-t)$ constant in $t$. Its limit is zero
by (GP.12) and boundedness of $A$, whereas some value is nonzero:
avoid the two discrete sets of zeros of $A(t)$ and $A(u-t)$.
For $k>2$, choose fixed $t_2,\ldots,t_{k-1}$ at nonzeros of $A$
and vary $t_1=t$, $t_k=u-t-\sum_{2\leq i<k}t_i$. The identical
argument makes the product both zero and nonzero, a contradiction.
The assertion for $\gamma_k$ follows by dividing by the strictly
positive $b_k$ already proved. This proof uses only the explicit
quartet tail, not a conjecture about positions of unselected zeros.
\end{proof}

\section{Full complex coefficient transform and exact variance norm}
Put $\alpha=2\lambda$ and keep the original monic polynomials
\begin{equation}
 \sum_{n\geq0}\frac{b_n^{(\lambda)}(t)}{n!}z^n
 =(1-iz)^{-\lambda+it/2}(1+iz)^{-\lambda-it/2},\qquad
 \|b_n^{(\lambda)}\|^2=c_\lambda n!(\alpha)_n.
 \tag{GP.14}
\end{equation}
The specified gamma measure has exponential moments. Its polynomials
are complete in $L^2$: for a vector orthogonal to them, multiplying
it by the measure gives a finite complex measure by Cauchy--Schwarz;
the same inequality with exponential moments makes its Fourier
transform holomorphic on a strip about the real axis. All derivatives
at zero vanish. Analytic uniqueness and Fourier uniqueness make the
measure, hence the vector, zero. The same proof applies to the sum
measure with mass $c_\lambda^k$.

Define complex coefficients, retaining conjugation conventions,
\begin{equation}
 \mathfrak a_j=\frac{\int b_j^{(\lambda)}(t)A(t)r_\lambda(t)dt}
                         {c_\lambda j!(\alpha)_j},\quad
 \mathfrak A(z)=\sum_{j\geq0}(\alpha)_j\mathfrak a_jz^j,\quad
 e_{k,n}=[z^n]\mathfrak A(z)^k.
 \tag{GP.15}
\end{equation}
The square-integrable vector $A$ has these orthogonal coefficients.
The series in (GP.15) may be used coefficientwise; no interchange of
an infinite product with an integral is required below.
\begin{theorem}[Complex convolution and its exact lost norm]
With all original masses,
\begin{equation}
 \begin{gathered}
 a_k(u)=\sum_{n\geq0}\frac{e_{k,n}}{(k\alpha)_n}b_n^{(k\lambda)}(u)
       \quad\hbox{in }L^2(r_{\lambda,k}du),\\
 \int d_k r_{\lambda,k}
  =\mu_h^k-c_\lambda^k\sum_{n\geq0}
                         \frac{n!}{(k\alpha)_n}|e_{k,n}|^2.
 \end{gathered}
 \tag{GP.16}
\end{equation}
For every polynomial $F$ in the original $S$ coordinate and every
real $|\theta|<\pi/2$, the exact weighted version is
\begin{equation}
 \begin{aligned}
 \|a\,U(e^{\theta u/2}F(k/2+iu))\|^2
 &=\int e^{\theta u}|F(k/2+iu)|^2|a_k(u)|^2r_{\lambda,k}(u)du\\
 &\quad+\int e^{\theta u}|F(k/2+iu)|^2d_k(u)r_{\lambda,k}(u)du.
 \end{aligned}
 \tag{GP.17}
\end{equation}
\end{theorem}
\begin{proof}
Multiplication of the finite-degree coefficients of the generating
functions gives
\[
 b_n^{(k\lambda)}(\textstyle\sum_i t_i)
 =\sum_{n_1+\cdots+n_k=n}\frac{n!}{\prod_i n_i!}
                                \prod_i b_{n_i}^{(\lambda)}(t_i).
\]
Pairing this finite identity with $a$ gives
$c_\lambda^kn!e_{k,n}$. Division by the original sum norm
$c_\lambda^kn!(k\alpha)_n$, followed by completeness, gives the
first identity. Pythagoras for $P=UC$ gives the second since
$\|a\|^2=\mu_h^k$. Finally conditional expectation commutes with
multiplication by a function of the sum. Apply the same orthogonal
decomposition to $aU(e^{\theta u/2}F)$ to obtain (GP.17).
Boundedness by $C_h^k$, (GP.2), and polynomial growth prove all
integrals finite, including any fixed number of theta derivatives.
\end{proof}

\section{The exact induced change of the original finite quotient metric}
Fix the actual annihilator $\chi=\chi_{h,k}$, degree $q>0$, and
$N\geq q-1$. Use original polynomial spaces
$E=\PP_N(S)$, $E_\partial=\PP_{N-q}(S)$, where a negative degree
means the zero space, and $C_\chi=\C[S]/(\chi)$. Let
$B_\chi:E_\partial\to E$ be multiplication by $\chi$ and
$\pi:E\to C_\chi$ the quotient. Use the original monomial bases,
or any explicitly transported fixed bases, in every matrix below.
For a coefficient column $x$ denote its polynomial by $F_x$.
Define three source matrices by
\begin{align}
 x^*M_hx&=\int e^{\theta u}|F_x(k/2+iu)|^2b_k r_{\lambda,k}du,
 \nonumber\\
 x^*M_{\rm coh}x&=\int e^{\theta u}|F_x(k/2+iu)|^2|a_k|^2r_{\lambda,k}du,
 \nonumber\\
 x^*M_{\rm rel}x&=\int e^{\theta u}|F_x(k/2+iu)|^2d_kr_{\lambda,k}du.
 \tag{GP.18}
\end{align}
These matrices are elements of $\operatorname{Herm}(E)$, equivalently
linear maps $E\to E^\vee$ with $E^\vee$ the conjugate dual under
$x\mapsto[y\mapsto y^*Mx]$. The associated linear observation maps are
$F\mapsto aU(e^{\theta u/2}F)$,
$F\mapsto U(e^{\theta u/2}a_kF)$, and
$F\mapsto (I-P)aU(e^{\theta u/2}F)$.
The Gram operation itself has type
$\operatorname{Hom}_{\C}(E,\HH)\to\operatorname{Herm}(E)$,
$R\mapsto R^*R$. Thus
\begin{equation}
 M_h=M_{\rm coh}+M_{\rm rel}.
 \tag{GP.19}
\end{equation}
The original arithmetic matrix $M_h$ is unchanged.

\begin{lemma}
$M_{\rm coh}$ is positive definite for every packet/reference admitted
above, every $k\geq1$, every finite $N$, and real $|\theta|<\pi/2$.
For the quartets of (GP.12) and $k\geq2$, $M_{\rm rel}$ is also
positive definite.
\end{lemma}
\begin{proof}
Put $f=Ar_\lambda$. It is a nonzero $L^1$ function, since $A$ is
bounded and nonzero almost everywhere. The numerator of $a_k$ is
$f^{*k}$. If this convolution were zero almost everywhere, its
Fourier transform $(\widehat f)^k$ would vanish everywhere, whence
$\widehat f=0$ and $f=0$ by Fourier uniqueness, a contradiction.
Consequently $|a_k|^2r_{\lambda,k}$ is positive on a set of positive
Lebesgue measure. A nonzero polynomial on the indicated affine line
has only finitely many real zeros, so its squared norm is strictly
positive. This proves the first assertion. The second follows
from $d_k(u)>0$ for every $u$, proved in (GP.13).
\end{proof}

Let $R_h,R_{\rm coh}:C_\chi\to E$ be the least-norm sections for
$M_h,M_{\rm coh}$, and write $G_h=R_h^*M_hR_h$,
$G_{\rm coh}=R_{\rm coh}^*M_{\rm coh}R_{\rm coh}$. They exist uniquely:
for any section $R_0$, subtract
$B_\chi(B_\chi^*MB_\chi)^{-1}B_\chi^*MR_0$; the boundary Gram
is invertible because $B_\chi$ is injective and $M$ positive definite.
An empty boundary contributes the zero map.
\begin{theorem}[Exact two positive terms in the quotient change]
In these original coordinates put
\begin{equation}
 K=(B_\chi^*M_hB_\chi)^{-1}B_\chi^*M_{\rm rel}R_{\rm coh}.
 \tag{GP.20}
\end{equation}
It is a map $C_\chi\to E_\partial$, or the unique zero map when
the boundary is empty. Then
\begin{equation}
 R_h=R_{\rm coh}-B_\chi K,\qquad
 \boxed{G_h=G_{\rm coh}
       +K^*B_\chi^*M_{\rm coh}B_\chi K
       +R_h^*M_{\rm rel}R_h.}
 \tag{GP.21}
\end{equation}
For an actual full quartet, $k\geq2$, and $q>0$, this gives
$G_h-G_{\rm coh}>0$ for every admitted $N$ and theta. More precisely,
with the actual positive matrix
\begin{equation}
 E_q=G_{\rm coh}^{-1/2}
 (K^*B_\chi^*M_{\rm coh}B_\chi K+R_h^*M_{\rm rel}R_h)
 G_{\rm coh}^{-1/2},
 \tag{GP.22}
\end{equation}
the original determinants satisfy
\begin{equation}
 \det G_h=\det G_{\rm coh}\det(I+E_q).
 \tag{GP.23}
\end{equation}
These comparison conjugations do not replace either quotient metric.
\end{theorem}
\begin{proof}
$B_\chi^*M_{\rm coh}R_{\rm coh}=0$, so the formula for the
$M_h$ section using $R_{\rm coh}$ as initial section gives (GP.20)
and the first identity. Expand
$R_h^*M_{\rm coh}R_h$ about $R_{\rm coh}$. The two cross terms
vanish by the stated orthogonality. Add $R_h^*M_{\rm rel}R_h$
using (GP.19); this proves (GP.21) with both terms retained.
$R_h$ is injective because $\pi R_h=I$. Thus strict positivity
of $M_{\rm rel}$ makes its displayed pullback positive definite
on the nonzero quotient. Conjugation by the actual positive square
root of $G_{\rm coh}$ proves (GP.23). For $N=q-1$ the first
positive term is zero and the second remains as stated. For $q=0$
the original quotient is zero: both section maps are zero, and the
empty determinant is one. The nonzero analytic source does not
supply an extra quotient dimension.
\end{proof}

\section{Support, the original unit, and the domain of arithmetic observation}
The original cyclic source is
$\V P=P(D_1+\cdots+D_k)F_h^{\otimes k}$. In its Mellin realization
its full image is
\begin{equation}
 \prod_i\frac{\Gamma(\lambda+it_i/2)}{\sqrt{2\pi}}
                  a(\mathbf t)P(k/2+i\textstyle\sum_i t_i).
 \tag{GP.24}
\end{equation}
Multiplication by the displayed gamma factors is an isometry
$\HH\to L^2(d\mathbf t)$. Equations (GP.9) and (GP.17) split
the vector before this multiplication and reconstruct it by addition.
On the image graph
\[
 \mathcal G_N=\{(U(a_kF), (I-P)aUF):F\in\PP_N(S)\}
\]
addition $\operatorname{add}_N:\mathcal G_N\to aU(\PP_N)$ is
injective, since the squared norm of the reconstructed vector is the
strictly positive polynomial norm $M_h$. It is onto the displayed
image by construction. Denote the multiplication in (GP.24) by
$M_\Gamma:\HH\to L^2(d\mathbf t)$ and the original unitary Mellin
map by $\mathcal M_{\rm unit}$. The graph observation is exactly
$J^{(k)}\mathcal M_{\rm unit}^{-1}M_\Gamma\operatorname{add}_N$,
and likewise with $q^{(k)}$. These are well-defined linear maps
because the reconstructed image is precisely the original admitted
Mellin source. Their values on the graph of $P$ are exactly
\begin{equation}
 J^{(k)}\V P=\eta_{h,k}\pi P,\qquad
 q^{(k)}\V P=\sigma_h^{\otimes k}\eta_{h,k}\pi P,
 \quad\eta_{h,k}[P]=\upsilon_h^{\otimes k}P(A_k)1.
 \tag{GP.25}
\end{equation}
This specifies the domain of the observation; no arithmetic jet is
assigned to an arbitrary independent measurable relative vector.
The graph in this display uses theta zero. At a nonzero admitted theta,
multiply both graph components by the same $e^{\theta u/2}$; its inverse
multiplies by $e^{-\theta u/2}$ before the original inverse Mellin map.
Both maps are isometries between the declared tilted and untilted spaces.

Here is the full cochain map for the relation. In the original tensor algebra
\[
 \C[s_1,\ldots,s_k]/(h(s_1),\ldots,h(s_k)),
\]
the cyclic annihilator satisfies $\chi(s_1+\cdots+s_k)=0$. Successive monic
division in the fixed order $s_1,\ldots,s_k$ gives unique quotients
at each division step and polynomials $Q_i$ with
\[
 \chi(s_1+\cdots+s_k)=\sum_{i=1}^k h(s_i)Q_i(s_1,\ldots,s_k).
\]
The final remainder is zero because the monomials of degree less
than $\deg h$ in each variable are a basis of that tensor algebra.
In $[V\xrightarrow{\Theta}\mathscr B]^{\otimes k}$ put
$D=-x\partial_x$ and
$\phi_0(x)=(4\pi^2x^4-6\pi x^2)e^{-\pi x^2}$. It is even,
$\phi_0(0)=0$ and $\int_\R\phi_0=0$: the two even Gaussian
integrals are $4\pi^2\cdot3/(4\pi^2)$ and
$6\pi\cdot1/(2\pi)$, both equal to three. For $\Re s>1$,
termwise Mellin integration of
$\Theta\phi_0(x)=\sum_{n\ne0}\phi_0(nx)$ yields
$2\zeta(s)\mathcal M\phi_0(s)$. Direct integration with
$y=\pi x^2$ gives
\[
 \mathcal M\phi_0(s)=\pi^{-s/2}
 \bigl(2\Gamma(s/2+2)-3\Gamma(s/2+1)\bigr)
 =\tfrac12s(s-1)\pi^{-s/2}\Gamma(s/2).
\]
Thus $\mathcal M(\Theta\phi_0)=g$ there and by continuation on
the original Mellin domain. The original inverse Mellin definition
of $F_h$ then gives $h(D)F_h=\Theta\phi_0$. The precise primitive is
\begin{equation}
 \mathcal P_\chi F=
 \sum_{i=1}^k(-1)^{i-1}Q_i(D_1,\ldots,D_k)
 F(D_1+\cdots+D_k)
 (F_h^{\otimes(i-1)}\otimes\phi_0\otimes F_h^{\otimes(k-i)}).
 \tag{GP.26}
\end{equation}
Euler differentiation preserves the even Schwartz source and its
two conditions: its value at zero vanishes and its integral is the
original integral by integration by parts. It also preserves rapid
decay at both multiplicative ends and commutes with $\Theta$.
Every term therefore lies in the asserted cochain domain. The
preceding $i-1$ degree-one factors give the tensor differential sign
$(-1)^{i-1}$, cancelling the written sign. Hence, with no sign omitted,
\begin{equation}
 d\mathcal P_\chi F=\V(\chi F),\qquad
 d[-\mathcal P_\chi K]=\V(R_h-R_{\rm coh}).
 \tag{GP.27}
\end{equation}
This proves the cochain correction map on $C_\chi$ in (GP.21),
including its zero value when the boundary is empty.

Every linear map on this graph and every original source or quotient
has its original fixed-support lift: $(\ell,x)\mapsto(\ell,Lx)$,
and external absence maps to external absence. The maps preserve
addition and scalar action because the underlying maps are linear.
For an actual relation the quotient value is $(\ell,0)$, its supported
zero, with the source and the two graph components still reconstructible.
The norm has type $\HH\to\R_{\geq0}$ and is positively homogeneous;
its square is quadratic. Trace has type
$\operatorname{End}_{\C}(E)\to\C$ and is linear; its restriction to
Hermitian endomorphisms is real-linear. The Gram map and squared
modulus are quadratic on the underlying real vector spaces, and
trace composed with a Gram map is quadratic. The exact linear
observation maps producing the source Grams and the quotient values
are (GP.18), (GP.21), and (GP.24)--(GP.25).

\section{Exact polynomial calibration with the original gamma mass}
The following auxiliary calibration tests the complete complex-phase
identities by rational arithmetic. It is a declared polynomial amplitude,
with all fixed moments finite in the original gamma measure. It does not
invoke the bounded packet estimate (GP.12). Take $\lambda=1$, $k=2$,
$c_1=1/2$, $c_1^2=1/4$, and $A(t)=1+it/2$. The original sum coordinate
is $S=1+iu$. Direct conditional polynomial calculation gives
\begin{equation}
 \begin{aligned}
 a_2(u)&=\frac65+\frac{iu}{2}-\frac{u^2}{20},\\
 b_2(u)&=\frac{54}{35}+\frac{39u^2}{280}+\frac{3u^4}{1120},\\
 d_2(u)&=\frac{18}{175}+\frac{13u^2}{1400}+\frac{u^4}{5600}>0.
 \end{aligned}
 \tag{GP.28}
\end{equation}
Here are enough exact polynomial identities to prove the calculation.
For the one-factor parameter $\alpha=2$, $b_2^{(1)}(t)=t^2-2$.
For the sum parameter $4$, the second and fourth monic polynomials are
$u^2-4$ and $u^4-32u^2+72$. The conditional addition formula yields
\begin{align*}
 C(t_1t_2)&=(u^2-4)/5,\\
 C(t_1^2)=C(t_2^2)&=3(u^2-4)/10+2,\\
 C(t_1^2t_2^2)&=\frac3{70}(u^4-32u^2+72)
                        +\frac65(u^2-4)+4.
\end{align*}
Indeed the respective coefficients are
$2^2/(4)_2$, $(2)_2/(4)_2$, and $(2)_2^2/(4)_4$.
Expanding $(1+it_1/2)(1+it_2/2)$ and
$(1+t_1^2/4)(1+t_2^2/4)$ using these identities gives the first two
lines of (GP.28). Subtracting the squared modulus gives its last line.
This is the complete complex subtraction, including its imaginary
linear coefficient.

For the auxiliary annihilator $\chi(S)=(S-1)^2-2$, use the original
$S$ monomials at $N=1$, where the boundary is empty. The sum gamma
moments of orders zero, two, four, six are respectively
$1/4,1,14,376$, obtained by differentiating
$\tfrac14\cos^{-4}\theta$. Substitution in (GP.18) gives
\begin{equation}
 \begin{gathered}
 G_h=\begin{pmatrix}9/16&9/16\\9/16&81/16\end{pmatrix},\quad
 G_{\rm coh}=\begin{pmatrix}21/40&21/40\\21/40&189/40\end{pmatrix},\\
 G_h-G_{\rm coh}=\begin{pmatrix}3/80&3/80\\3/80&27/80\end{pmatrix},\\
 \det G_h=81/32,\quad\det G_{\rm coh}=441/200,\quad
 \det(G_h-G_{\rm coh})=9/800.
 \end{gathered}
 \tag{GP.29}
\end{equation}
For example $\int b_2r_{1,2}=9/16$ and
$\int u^2b_2r_{1,2}=9/2$, while the corresponding coherent moments
are $21/40$ and $21/5$. Odd moments vanish, so the original
off-diagonal $S$ entries equal the masses. These calculations prove
every entry above. They exhibit the strictly positive contribution
in the unchanged quotient coordinates.

The accompanying checker additionally derives moments from the Laplace
function, solves the independent product-moment equations, and checks
all moments through degree eight and the section/quotient identities
through $N=4$. It uses a separate two-atom conditional fibre of original
masses two and three to test the complex adjoint and both projection
traces. Its mathematical controls change a conjugate transpose, omit
an original mass factor, or omit the relative quotient term. Those
finite checks supplement the proofs above and do not estimate an
actual growing-degree zeta quotient.

\section{Provenance and consequence}
The source constants, packet, conditional density, and coefficient
addition identity come from the complete delivered Gamma Convolution
Descent note, equations (2), (5)--(19), (24)--(27), and (32).
The original unit and theta relation primitive come from the retained
cyclic source construction and TVB in the cumulative reader. The
classical generating function is
\href{https://dlmf.nist.gov/18.23.E7}{NIST DLMF 18.23.7}, with
$b_n^{(\lambda)}(t)=n!P_n^{(\lambda)}(t/2;\pi/2)$.
The exact specialization and norms are proved in TG.7--TG.14.
The bounded amplitude is used here to prove the full phase transport,
fibre projection formulas, strict quartet variance, complex coefficient
transform, and two-term positive quotient identity. These equations
provide additional exact data for the same Toda quotient. They do not
give a growing-degree estimate for its consecutive volume ratio.
That ratio remains the actual determinant expression of the retained
metric, now with the explicit contribution (GP.21)--(GP.23).
\end{document}
